\subsection*{5.4 Proof of Lemma 21}
To prove Lemma21 it is sufficient to prove:
\begin{authorclaim}{Proposition 36}
There exist matrices $T_1,\ldots,T_k\in\mathbb F_2^{2^k\times k}$ such that (1) for $\alpha_i\in\{0,1\}$, $\sum_i\alpha_iT_i=0$ iff all $\alpha_i=0$, i.e. they are linearly independent; (2) every nonzero $M\in\operatorname{Span}(T_1,\ldots,T_k)$ satisfies $M+K_k\in S_k$.
\end{authorclaim}
\begin{proof}[Proof of Lemma21]
Let for $i\in\{1,\ldots,k-1\}$, $A_i(0)=T_i$ and $A_i(1)$ be the zero matrix. Let $A_k(0)=K_k+T_k$, $A_k(1)=K_k$. For each $b_1,\ldots,b_k\in\{0,1\}$, $\sum_{i=1}^kA_i(b_i)=\sum_{i=1}^k(1-b_i)T_i+K_k$. Then $\sum_{i=1}^kA_i(1)=K_k$, and if for at least one $i\in[k]$, $b_i\ne1$, then by the first condition of Proposition36, $\sum_{i=1}^k(1-b_i)T_i$ differs from zero, thus by the second condition of Proposition36, $\sum_{i=1}^kA_i(b_i)\in S_k$.
\end{proof}
\begin{proof}[Proof of Proposition36]
Let us prove the proposition by induction on $k$. We are going to prove a stronger statement: namely, we additionally require that for arbitrary non-zero matrix $M\in\operatorname{Span}(T_1,\ldots,T_k)$ the set of even-indexed rows of $M+K_k\in S_k$ coincide with the set of odd-indexed rows of this matrix with all bits flipped.
The base case $k=1$. $T_1=\binom01$, and $K_1=\binom00$. It is easy to verify that all conditions hold.
Induction step from $k$ to $k+1$. Notice that
\[
K_{k+1}=\begin{pmatrix}K_k&0_{2^k\times1}\\K_k&1_{2^k\times1}\end{pmatrix}.
\]
Let $T_1,\ldots,T_k$ be the matrices from induction hypothesis for $k$. Then define
\[
T'_i=\begin{pmatrix}T_i&0_{2^k\times1}\\T_i&0_{2^k\times1}\end{pmatrix}
\]
for $i\in[k]$ and
\[
T'_{k+1}=\begin{pmatrix}0_{2^k\times k}&z_0\\1_{2^k\times k}&z_1\end{pmatrix},
\]
where $z_0=(0,1,0,1,\ldots,0,1)^T\in\{0,1\}^{2^k\times1}$, and $z_1=(1,0,1,0,\ldots,1,0)^T\in\{0,1\}^{2^k\times1}$.
Let us verify that all conditions hold. First we show that the matrices $T'_1,T'_2,\ldots,T'_{k+1}$ are linearly independent. Matrices $T'_1,T'_2,\ldots,T'_k$ are linearly independent since they contain linearly independent blocks $T_1,\ldots,T_k$. The matrix $T'_{k+1}$ does not belong to $\operatorname{Span}(T'_1,\ldots,T'_k)$, since the last column of $T'_{k+1}$ is non-zero, but the last columns of all $T'_1,\ldots,T'_k$ are zeros.
Let us check that for any non-zero matrix $M\in\operatorname{Span}(T'_1,\ldots,T'_k,T'_{k+1})$, the condition $M+K_{k+1}\in S_{k+1}$ holds and the set of even-indexed rows of $M+K_{k+1}$ coincide with the set of odd-indexed rows of this matrix with all bits flipped. Let us analyze the cases:
\begin{enumerate}
\item Let $M$ be a non-zero matrix from $\operatorname{Span}(T'_1,\ldots,T'_k)$. Then, $M$ has form
\[
\begin{pmatrix}M'&0_{2^k\times1}\\M'&0_{2^k\times1}\end{pmatrix},
\]
where $M'$ is a non-zero matrix from $\operatorname{Span}(T_1,\ldots,T_k)$, thus $M'+K_k\in S_k$. Then
\[
M+K_{k+1}=\begin{pmatrix}M'+K_k&0_{2^k\times1}\\M'+K_k&1_{2^k\times1}\end{pmatrix};
\]
it follows from the induction hypothesis that all rows of this matrix are distinct, i.e. $M+K_{k+1}\in S_{k+1}$. In order to verify that the set of even-indexed rows of this matrix coincide with the set of odd-indexed rows with all bits flipped, observe that by induction hypothesis the first $2^{k-1}$ even-indexed rows of $M+K_{k+1}$ coincide with the last $2^{k-1}$ odd-indexed rows of $M+K_{k+1}$ with all bits flipped, and the first $2^{k-1}$ odd-indexed rows of $M+K_{k+1}$ coincide with the last $2^{k-1}$ even-indexed rows of $M+K_{k+1}$ with flipped bits.
\item $M=T'_{k+1}$, then
\[
M+K_{k+1}=\begin{pmatrix}K_k&z_0\\1_{2^k\times k}+K_k&z_0\end{pmatrix}.
\]
Let us show that all rows of this matrix are distinct. The first $2^k$ rows start with0 and are obtained by appending zeroes and ones to the rows of $K_k$ in the alternating order. Since for every pair of coinciding rows of $K_k$ they are adjacent, the first $2^k$ rows are distinct. The last $2^k$ rows start from one, so they differ from the first $2^k$ rows. The proof that they are distinct is the same as for the first $2^k$ rows. Observe that the $(2i-1)$th row of the matrix $M+K_{k+1}$ coincide with the $(2^k+2i)$th row of $M+K_{k+1}$ with flipped bits, and the $(2i)$th row of $M+K_{k+1}$ coincide with the $(2^k+2i-1)$th row of $M+K_{k+1}$ with flipped bits for $i\in[2^k]$.
\item $M=R+T'_{k+1}$, where $R$ is a non-zero matrix from $\operatorname{Span}(T'_1,\ldots,T'_k)$. Let $R$ have the form
\[
\begin{pmatrix}R'&0_{2^k\times1}\\R'&0_{2^k\times1}\end{pmatrix},
\]
where $R'$ is a non-zero matrix from $\operatorname{Span}(T_1,\ldots,T_k)$. Then
\[
M+K_{k+1}=R+T'_{k+1}+K_{k+1}
=\begin{pmatrix}R'+K_k&z_0\\1_{2^k\times k}+R'+K_k&z_0\end{pmatrix}.
\]
By the induction hypothesis, $R'+K_k\in S_k$ and its even-indexed rows coincide with its odd-indexed rows with flipped bits. Then, all even-indexed rows of $M+K_{k+1}$ end with0, the first $2^{k-1}$ of them are even-indexed rows of $R'+K_k$ with appended zero, and the last $2^{k-1}$ of them are even-indexed rows of $R'+K_k$ with all bits flipped and appended0. Then, by the induction hypothesis, the set of the former rows does not intersect with the set of the latter rows, therefore they are all distinct. By the same argument, all the rows of $M+K_{k+1}$ that end with1 are distinct. Thus, $M+K_{k+1}\in S_{k+1}$.
Let us verify that the set of even-indexed rows of this matrix coincide with the set of odd-indexed rows of this matrix with all bits flipped. Observe that if the $i$th row of $R'+K_k$ coincides with the $j$th row of $R'+K_k$ with flipped bits, then the $i$th row of $M+K_{k+1}$ coincides with its $j$th row with flipped bits, and the $(2^k+i)$th row of $M+K_{k+1}$ coincides with its $(2^k+j)$th row with all bits flipped. The required property follows from the induction hypothesis.
\end{enumerate}
\end{proof}
\subsection*{5.5 Corollaries}
\begin{authorclaim}{Corollary37}
If $k+7\le\ell$, the size of any semantic $\mathrm{Res}(\mathrm{PC}_{k-1})$ tree-like refutation of $\mathrm{BPHP}_{2^\ell}^{2^\ell+2^k}$ is at least $2^{\Omega(2^{\ell/2}/(k2^{3k/2}))}$. For $k=2$, every tree-like semantic $\mathrm{Res}(\oplus)$ refutation of $\mathrm{BPHP}_{2^\ell}^{2^\ell+4}$ has size at least $2^{\Omega(2^\ell)}$.
\end{authorclaim}
\begin{proof}
Follows from Theorem3 and Lemma1.
\end{proof}
\begin{authorclaim}{Corollary38}
Let $2\le k\le\ell-7$ and $S$ be the minimal size of tree-like refutation of $\phi=\mathrm{BPHP}_{2^\ell}^{2^\ell+2^k}\circ\oplus_k$ in the semantic proof system $\mathrm T^{\mathrm{cc}}(k,c)$. Then
$\log S\log\log S\ge c\cdot\Omega(2^{\ell/2}/(k2^{3k/2}))$.
For $k=2$, $\log S\log\log S\ge c\cdot\Omega(2^\ell)$.
\end{authorclaim}
\begin{proof}
By Lemma9, $R^{1/3}_{\mathrm{pub}}(\operatorname{Search}(\phi))=O(\log S\log\log S/c)$. Also
\[
R^{1/3}_{\mathrm{pub}}(\operatorname{Search}(\mathrm{BPHP}_{2^\ell}^{2^\ell+2^k}\circ\oplus_k))
\ge R^{1/3}_{\mathrm{pub}}(\oplus_k\operatorname{Search}(\mathrm{BPHP}_{2^\ell}^{2^\ell+2^k})).
\]
Now apply Theorem3.
\end{proof}
